\documentclass[11pt]{article}

\usepackage[margin=1.25in]{geometry}
\usepackage{amsmath}
\usepackage{amssymb}
\usepackage{booktabs}
\usepackage{longtable}
\usepackage{array}
\usepackage{float}
\usepackage{hyperref}
\usepackage{xcolor}
\usepackage{caption}
\usepackage{natbib}
\hypersetup{colorlinks=true, linkcolor=blue, citecolor=blue, urlcolor=blue}
\setlength{\parskip}{6pt}
\setlength{\parindent}{0pt}
\renewcommand{\baselinestretch}{1.25}

\title{Pricing a Quarterfinal Under a Cold-Start Crowd: \\
Integrating StatsBomb Event-Level History into the JTC Forecasting Pipeline \\
for France versus Morocco (2026 FIFA World Cup)}
\author{Souparneya Chakrabarti}
\date{July 9, 2026}

\begin{document}
\maketitle

\begin{abstract}
This paper documents the pricing of fifteen probability-market questions for the France
versus Morocco quarterfinal of the 2026 FIFA World Cup, submitted to the Jump Probability
Cup (JTC) on \textit{play.sportspredict.com}. The exercise differs from the forty-nine
prior matches priced in this campaign in two respects. First, no crowd consensus or sharp
market price was available for any of the fifteen questions at pricing time, so the
campaign's usual anchor---an average of crowd consensus and a verified sharp market---could
not be applied to any question; every estimate had to be built from first principles.
Second, this is the first match priced with access to a newly constructed StatsBomb
event-level panel spanning the 2018 and 2022 World Cups (128 matches, 6{,}131
player-match rows, 257 team-match rows), including the one prior competitive meeting
between these two sides, the 2022 semifinal (France 2--0 Morocco). The paper's central
methodological concern is stated directly by the research brief that motivated it: with
a much richer data surface available than in any previous match, the risk shifts from
data scarcity to indiscriminate application---fitting a tool to data that happens to
exist rather than data that is actually informative for the question at hand. Section
\ref{sec:data} inventories every locally available data source and classifies each as
relevant or irrelevant to this specific match. Section \ref{sec:methods} re-derives the
project's core statistical machinery---a point-in-time Elo rating system, a Poisson
goal-rate regression with Dixon--Coles low-score correction and Negative-Binomial
overdispersion, and an ordered logit outcome model---and reports the coefficients each
model contributes to this match. Section \ref{sec:questionmapping} maps each of the
fifteen questions to the specific data source and tool judged appropriate to it, with an
explicit rationale for rejecting the StatsBomb-only or model-only alternative in each
case where prior campaign evidence recommends blending. Section \ref{sec:results}
reports the fifteen final prices. Section \ref{sec:discussion} evaluates the exercise
against the campaign's own accumulated record of successes and catastrophic failures,
in particular the finding that projecting a team shot-on-target or corner count upward
from a thin empirical sample---absent a market to anchor to---has been the single most
expensive repeated error of the campaign (net $-155.64$ relative Brier points across
five prior matches), a risk this match faces acutely on two of its fifteen questions
given the complete absence of a market anchor. Section \ref{sec:marketrevision}
documents a subsequent revision of these estimates after live external market data was
retrieved. An initial check of the JTC platform's own web API was misread as showing a
still-unformed crowd consensus; a follow-up validation against markets with known trade
activity and known settlement outcomes shows this reading was incorrect, and that the
field in question is a binary open/settled flag rather than a probability at all---this
project's pre-existing finding that the pre-close crowd consensus is simply not
observable via any known API endpoint stands uncorrected. A public Smarkets sharp-market
query, unaffected by this issue, supplied a real, liquid price for twelve of the fifteen
questions, most consequentially raising the match-winner estimate from an
ordered-logit-only 0.52 to a market-anchored 0.60.
\end{abstract}

\section{Introduction}

The Jump Probability Cup (JTC) is a season-long forecasting tournament in which
participants submit a probability for each of a fixed set of yes/no questions attached to
a football match, and are scored against the platform's own aggregate crowd consensus
using a Relative Brier Points (RBP) metric: positive when the submission beats the crowd
on a given question, negative when it does not. Because RBP is scored relative to the
crowd rather than against an absolute outcome, the object of interest throughout this
campaign has never been calibration in the abstract but calibration relative to a
specific, empirically compressible crowd. Prior work in this campaign (documented in
\texttt{CROWD\_BIAS\_REGRESSION\_UPDATE.md}) established that the JTC crowd systematically
compresses extreme probabilities toward one half: fit on 384 previously scored group-stage
questions, $\hat p_{\text{crowd}} \approx 0.511 + 0.600\,(\hat p_{\text{model}} - 0.5)$,
with a 95\% confidence interval on the slope of $[0.559, 0.641]$ that excludes both zero
and one. Because the crowd under-moves relative to a well-calibrated model, a genuinely
informative model that is not itself compressed toward 50\% has a structural, repeatable
edge.

That structural edge, however, depends on having a crowd to compress against. The match
addressed in this paper---France versus Morocco, kickoff July 9, 2026, a Quarterfinal of
the 2026 World Cup---presented a condition not previously encountered at this stage of the
campaign: at pricing time, the question set had just opened and no crowd consensus, and no
Smarkets or other sharp market quote, existed for any of the fifteen questions. A
repository-wide search of every location where market or crowd data is recorded in this
project (\texttt{bot/data/markets\_raw.jsonl}, \texttt{bot/data/matches.json},
\texttt{data/external\_markets/}, and the \texttt{matches/} directory) confirmed that all
fifteen markets carried a bare \texttt{status: open} field with no accompanying price. The
campaign's foundational rule for match-winner pricing---submit the average of the crowd
consensus and a verified sharp market (\texttt{RULE1})---and its counterpart for
props---submit at the market price when a market exists (\texttt{RULE2})---were therefore
inapplicable to every one of the fifteen questions. This paper is, in effect, a controlled
test of what the project's statistical and heuristic apparatus can produce with no live
external anchor at all, leaning instead on point-in-time historical models and a newly
available body of event-level data.

The second novel feature of this match is data abundance rather than data scarcity. Over
the preceding several days the project constructed, validated, and documented (in
\texttt{STATSBOMB\_DATASET\_AUDIT.md} and
\texttt{STATSBOMB\_INTEGRATION\_AND\_STATS\_TESTS\_2026-07-08.md}) a flattened panel of
StatsBomb open-data event and lineup files covering the 2018 and 2022 World Cups, giving,
for the first time in this campaign, individual shot, shot-on-target, card, and minutes
data for named players across an entire major tournament history rather than only the
current tournament's four-or-five-game sample. This is a genuine increase in the
information available for this match: three of the four named players in the question
set (Mbapp\'e, Dembel\'e, Hakimi) have ten-to-fourteen-match StatsBomb histories, and the
2022 semifinal between these exact two teams is itself in the panel. The methodological
risk this introduces is the one identified in the brief that motivated this paper: a
richer data surface increases the temptation to apply a sophisticated tool because the
data to feed it exists, rather than because the tool is the right one for the specific
question and because the resulting estimate has actually been validated against this
campaign's own record. Section \ref{sec:questionmapping} treats this concern explicitly,
and Section \ref{sec:discussion} reports a diagnostic already run on exactly this
question in \texttt{STATSBOMB\_INTEGRATION\_AND\_STATS\_TESTS\_2026-07-08.md}: a
pure StatsBomb historical base-rate model, run head-to-head against the campaign's actual
submission on a different quarterfinal-adjacent match (Portugal versus Croatia), scored
93.01 RBP-equivalent against the submission's actual 120.30---losing overall despite
individually beating six of fourteen comparable questions. The documented conclusion of
that diagnostic, adopted as governing practice here, is that StatsBomb history is best
used as a shrinkage prior underneath the current tournament's thin live sample, not as a
standalone replacement for it.

\section{Data Inventory and Reliability Classification}
\label{sec:data}

Before any model was fit or any question priced, every locally available data source was
inventoried and classified by whether it carries information specific to France, Morocco,
or the four named players in this question set (Kylian Mbapp\'e, Ousmane Dembel\'e,
Achraf Hakimi, Brahim D\'iaz), as opposed to being generically available but uninformative
for this particular match. Table \ref{tab:datainventory} summarizes the classification.

\begin{table}[H]
\centering
\caption{Data source inventory, classified by relevance to France vs. Morocco.}
\label{tab:datainventory}
\begin{tabular}{p{4.6cm}p{3.0cm}p{6.2cm}}
\toprule
Source & Classification & Note \\
\midrule
StatsBomb event/lineup JSON, WC2018+WC2022 (\texttt{data/external/statsbomb/open-data/}) &
Rich and relevant & Contains the exact 2022 France--Morocco semifinal (match\_id 3869552); full events, lineups, and 360$^\circ$ tracking. \\
StatsBomb player-match panel (\texttt{statsbomb\_player\_match\_panel.csv}) &
Rich and relevant, uneven & 14 rows for Mbapp\'e, 21 for Dembel\'e, 10 for Hakimi, 0 for Brahim D\'iaz (uncapped by Morocco during 2018/2022). \\
StatsBomb team-match panel (\texttt{statsbomb\_team\_match\_panel.csv}) &
Rich and relevant & 28 France rows, 20 Morocco rows across WC2018/2022; includes the 2022 h2h team-level line. \\
ESPN match/rolling panels (\texttt{espn\_match\_panel.csv}, \texttt{espn\_rolling\_averages.csv}) &
Rich and relevant & Full 2026 group-stage box scores for both teams; does not yet include the R32/R16 knockout matches (those live only in the \texttt{matches/} folder artifacts). \\
Elo panel and current ratings (\texttt{elo\_match\_panel.csv}, \texttt{elo\_current\_ratings.csv}) &
Rich and relevant & Point-in-time ratings through the most recent settled match for both teams; source of the Elo differential used throughout Section \ref{sec:methods}. \\
\texttt{matches/France\_vs\_Paraguay/}, \texttt{matches/Canada\_vs\_Morocco/} &
Rich and relevant & Each team's most recent completed match, full pipeline artifacts including settled outcomes for the same named players. \\
\texttt{data/international\_results/results.csv} &
Rich and relevant & Full France--Morocco head-to-head since 1988 (49{,}473 total historical matches); score-only, no shot/card granularity outside the StatsBomb-covered 2022 fixture. \\
Transfermarkt national-team roster extract &
Present but unreliable & Zero France rows and 45 Morocco rows, none of which include Hakimi or Brahim D\'iaz; flagged as a gap rather than used. \\
StatsBomb open-data outside WC2018/2022 (club competitions, other tournaments) &
Rich but not relevant & Excellent for generic base-rate construction but contributes no team- or player-specific signal for this match beyond what the WC2018/2022 subset already provides. \\
Historical odds (\texttt{data/external/odds/}) &
Historically relevant only & Contains the actual closing odds for the 2022 semifinal but no 2026 quotes for this fixture. \\
\bottomrule
\end{tabular}
\end{table}

The single most consequential negative finding in this inventory is the confirmed absence
of any crowd consensus or market quote for this match anywhere in the repository, discussed
in the Introduction. The single most consequential positive finding is that the 2022
World Cup semifinal between these exact two teams is fully covered by the StatsBomb
event data, giving a genuine (if small, $n=1$) head-to-head precedent at the shot and
shot-on-target level, not merely the scoreline: in that match Mbapp\'e recorded three shots
and zero shots on target in 95 minutes, Dembel\'e recorded zero shots, and Hakimi recorded
one shot and zero shots on target, against a Morocco defense that conceded only once across
the entire 2022 tournament before that game.

\section{Methodology}
\label{sec:methods}

The forecasting pipeline used in this campaign, and applied here, follows the structure
of \citet{dixoncoles1997} and \citet{karlisntzoufras2003} for the goal-scoring process,
supplemented by an Elo-based outcome model following \citet{hvattum2010}. Each component
is re-derived below with the coefficients actually fitted on this project's own historical
panel of 49{,}400 international matches (1872--present), and the specific inputs this
match supplies to each.

\subsection{Point-in-Time Elo Ratings}

Ratings are updated match-by-match using the standard World Football Elo formula,
\[
R_{\text{new}} = R_{\text{old}} + K \cdot G \cdot (W - W_e),
\]
where $K$ is drawn from a tournament-tier lookup table (60 for World Cup finals matches,
30 for most qualifiers, 20 for friendlies), $G$ scales with margin of victory
($G=1$ for a one-goal difference, rising to $(11+|\Delta|)/8$ for a three-or-more-goal
difference), $W \in \{1, 0.5, 0\}$ is the match result from the home team's perspective,
and
\[
W_e = \frac{1}{10^{-\Delta r / 400} + 1}, \qquad
\Delta r = \begin{cases} (R_{\text{home}} - R_{\text{away}}) + 100 & \text{if not neutral} \\ R_{\text{home}} - R_{\text{away}} & \text{if neutral.} \end{cases}
\]
Because the ratings are recomputed strictly forward through time with no re-fitting on
future results, they carry no lookahead bias. As of the most recent settled match for each
side, France's rating stands at 2129.05 and Morocco's at 1984.85, an Elo differential of
$+144.2$ in France's favor. Because the 2026 World Cup is co-hosted by the United States,
Mexico, and Canada, and neither France nor Morocco is a host nation, this fixture is
treated as a neutral-venue match: no home-advantage term is applied to either side in any
downstream model.

\subsection{Poisson Goal-Rate Regression}

Goals are modeled as a Poisson process with rate
\[
\log(\lambda) = b_0 + b_1 \cdot \mathrm{home\_adv} + b_2 \cdot \Delta\mathrm{Elo},
\]
fit by maximum likelihood with an exponential recency weight
($\xi = 0.0008\text{/day}$, an approximately 2.4-year half-life) on 49{,}400 matches
(98{,}800 team-match goal observations). Table \ref{tab:poissoncoefs} reports the fitted
coefficients.

\begin{table}[H]
\centering
\caption{Poisson goal-rate regression coefficients.}
\label{tab:poissoncoefs}
\begin{tabular}{lrr}
\toprule
Parameter & Estimate & Std.\ Error \\
\midrule
Intercept ($b_0$) & 0.1041 & --- \\
Home advantage ($b_1$) & 0.2302 & 0.0211 \\
Elo differential ($b_2$, per Elo point) & 0.001810 & 0.0000368 \\
\bottomrule
\end{tabular}
\end{table}

With $\mathrm{home\_adv} = 0$ for both sides (neutral venue) and $\Delta\mathrm{Elo} =
\pm 144.2$, this gives
\[
\lambda_{\text{France}} = \exp(0.1041 + 0.001810 \times 144.2) = 1.441, \qquad
\lambda_{\text{Morocco}} = \exp(0.1041 - 0.001810 \times 144.2) = 0.855,
\]
a combined expected total of 2.295 goals. This total is markedly below the average goal
output either team has actually produced this tournament (France 2.8 goals scored per
game across five matches; Morocco 2.0), a discrepancy attributable to the relative
weakness of the group-stage and early-knockout opposition both sides have faced so far
(Senegal, Iraq, Norway, Sweden, and Paraguay for France; Brazil, Scotland, Haiti, and the
Netherlands for Morocco) rather than to any deficiency in the model, since $\Delta$Elo
already reflects the quality of today's specific opponent rather than the average quality
of past opponents.

\subsection{Dixon--Coles Correction and Negative-Binomial Overdispersion}

Following \citet{dixoncoles1997}, the independent-Poisson joint distribution is corrected
at the four low-score cells $(0,0), (1,0), (0,1), (1,1)$ by a multiplicative factor
$\tau(\cdot)$ governed by a single dependence parameter $\rho$:
\[
\tau(0,0) = 1 - \lambda_h \lambda_a \rho, \quad
\tau(1,0) = 1 + \lambda_a \rho, \quad
\tau(0,1) = 1 + \lambda_h \rho, \quad
\tau(1,1) = 1 - \rho.
\]
The project's own maximum-likelihood fit on the Poisson marginals gives $\hat\rho =
-0.06$; a Negative-Binomial refit of the marginals (to address a separately documented
overdispersion pattern in the goals-based pipeline) gives $\hat\rho_{\text{NB}} = -0.05$
alongside a dispersion parameter $\hat\alpha = 0.0992$ (SE $= 0.0131$, $n=98{,}800$),
confirming mild but statistically real overdispersion relative to the pure Poisson case
($\alpha = 0$). Under the independent-Poisson total ($\lambda_{\text{total}} = 2.295$),
$P(\text{total goals} \le 2) = 0.597$; the Negative-Binomial correction reallocates a
small amount of mass into the tails, which combined with the observation that the sole
StatsBomb-verified prior meeting between these two sides finished 2--0 (itself a
$\le 2$-goal outcome), motivates the slightly higher final estimate reported in Section
\ref{sec:results}.

\subsection{Ordered Logit Match Outcome Model}

The Poisson/Dixon--Coles pipeline is known, from an expanding-window backtest across the
full historical panel, to systematically underpredict the win probability of the stronger
side by roughly five to eight percentage points at every decile of predicted home-win
probability---a compression bias whose leading explanation is attenuation from
measurement noise in the Elo covariate itself, though this has not been fully closed out
diagnostically. To avoid inheriting this bias, match-winner probabilities in this campaign
are instead taken from an ordered logit fit directly on historical Win/Draw/Loss outcomes,
following the Elo-based outcome regression approach of \citet{hvattum2010}:
\[
\eta = b_{\text{elo}} \cdot \Delta\mathrm{Elo} + b_{\text{home}} \cdot \mathrm{home\_adv},
\qquad
P(\text{away win}) = \sigma(c_1 - \eta), \quad
P(\text{draw}) = \sigma(c_2 - \eta) - \sigma(c_1 - \eta), \quad
P(\text{home win}) = 1 - \sigma(c_2 - \eta),
\]
with $\sigma(\cdot)$ the logistic function. Table \ref{tab:ologitcoefs} reports the fitted
coefficients (49{,}400 matches).

\begin{table}[H]
\centering
\caption{Ordered logit match-outcome coefficients.}
\label{tab:ologitcoefs}
\begin{tabular}{lrr}
\toprule
Parameter & Estimate & Std.\ Error \\
\midrule
Elo differential ($b_{\text{elo}}$, per point) & 0.005199 & 0.000176 \\
Home advantage ($b_{\text{home}}$) & 0.3771 & 0.0745 \\
Cutpoint 1 ($c_1$) & $-0.7702$ & --- \\
Cutpoint 2 ($c_2$) & 0.5549 & --- \\
\bottomrule
\end{tabular}
\end{table}

With $\Delta\mathrm{Elo} = 144.2$ and $\mathrm{home\_adv} = 0$, $\eta = 0.7497$, giving
$P(\text{France win}) = 0.549$, $P(\text{draw}) = 0.272$, $P(\text{Morocco win}) = 0.180$.
This raw model output is adjusted downward by two to three percentage points in the final
submission (Section \ref{sec:results}) to reflect Morocco's demonstrated big-match
resilience this tournament---forcing the Netherlands to extra time and penalties in the
Round of 32, and a 3--0 defeat of Canada in the Round of 16---a qualitative adjustment kept
deliberately small in light of a separate finding in this campaign
(\texttt{STRATEGIC\_MARGIN\_PUSH\_RESEARCH.md}) that deliberately pushing a submitted
probability past one's honestly calibrated belief is strictly negative in expectation under
RBP's proper-scoring structure, except in a narrow leaderboard-cutoff scenario that does
not apply here.

\subsection{StatsBomb Empirical-Bayes Shrinkage for Player and Team Props}

For the player- and team-count questions (shots on target, corners, cards) that the
goals-based and ordered-logit models do not directly address, the campaign's existing
shrinkage methodology (\texttt{ml/statsbomb\_baserate\_test.py}) is applied: a team's rate
for a given statistic is shrunk toward the WC2018/2022 tournament-wide mean using
\[
\hat\mu_{\text{team}} = \frac{n \cdot \bar x_{\text{team}} + k \cdot \bar x_{\text{global}}}{n + k}, \qquad k = 5 \text{ pseudo-matches (team level)},
\]
and a named player's rate is shrunk toward a fixed weak prior (rather than a tournament
mean, there being too few directly comparable players to construct one) with $k = 3$
pseudo-matches. As documented in
\texttt{STATSBOMB\_INTEGRATION\_AND\_STATS\_TESTS\_2026-07-08.md}, a live test of this
exact method against an actual settled match (Portugal versus Croatia) found that the pure
StatsBomb-history model loses in aggregate to a tournament-aware blended judgment (93.01
versus 120.30 RBP-equivalent, on fourteen comparable questions), even though it wins
individually on six of the fourteen. The governing rule adopted here, consistent with that
finding, is that every StatsBomb-derived rate in Section \ref{sec:questionmapping} is used
as a prior to be blended with and generally subordinated to this tournament's own live
five-game form for each team, never as a standalone replacement for it.

\section{Question-by-Question Data and Tool Mapping}
\label{sec:questionmapping}

Table \ref{tab:mapping} states, for each of the fifteen questions, the primary data source
and pricing tool selected, together with the specific reason a plausible alternative tool
was not used instead. The purpose of this table is to make the avoid-overfitting discipline
requested for this exercise auditable rather than implicit: for every question, at least
one data-rich alternative existed that was deliberately not applied, and the table states
why.

\begin{longtable}{p{1.6cm}p{4.6cm}p{3.0cm}p{5.0cm}}
\caption{Question-to-tool mapping and the alternative deliberately rejected.} \label{tab:mapping} \\
\toprule
Q & Question (abbreviated) & Tool used & Rejected alternative and reason \\
\midrule
\endfirsthead
\toprule
Q & Question (abbreviated) & Tool used & Rejected alternative and reason \\
\midrule
\endhead
1 & Mbapp\'e anytime goal & StatsBomb career rate (0.857/game) blended down for opponent-specific regime check & A flat current-tournament hot-streak rate (4 of 5 games) was rejected as the sole input because it has never faced a defense of Morocco's caliber this cycle; the one StatsBomb-covered meeting with this exact opponent shows suppression. \\
2 & Hakimi 1+ SOT & StatsBomb career rate (0/10 tracked matches) & A generic starting-fullback positional prior was rejected in favor of his own, unusually informative, zero-SOT ten-match record. \\
3 & Brahim D\'iaz 1+ SOT & Current-tournament personal rate (1 of 5 games) & A team-context or positional prior was rejected because his own five-game personal sample is directly on point and this exact suppression pattern already won twice this campaign. \\
4 & Dembel\'e 2+ SOT & Current-tournament bimodal-by-opponent-quality pattern & A StatsBomb career average was rejected because it does not capture the opponent-quality gating that a direct within-tournament comparison (blanked vs.\ tougher sides, including the most recent match) already demonstrates. \\
5 & Morocco 3+ SOT & Team-level empirical floor from competitive games only & Blending upward toward a market-implied ceiling was rejected outright since no market exists to blend toward, and this exact move (projecting a team SOT count upward from a thin sample) produced the single worst result of the campaign in a prior match. \\
6 & Total goals $\le 2$ & Poisson/Dixon--Coles/NB total-goal model & A pure current-tournament scoring-rate average was rejected as it reflects weak group-stage opposition, not two elite defenses facing each other. \\
7 & 5+ total cards & Team disciplinary rate, deliberately not taken at face value & The raw historical average (very low for both teams) was not used unadjusted, since elimination-stage matches carry a documented general uplift in foul/card risk that a group-stage average cannot see. \\
8 & Tied at halftime & Tournament-wide empirical base rate & A team-specific first-half scoring model was rejected for lack of a large enough within-tournament first-half sample; the tournament-wide rate is the more stable estimator here. \\
9 & Both halves same goals & Derived consistency check against Q6/Q8 & No standalone tool exists for this question type in the project; priced for consistency with the two related direct questions rather than independently, to satisfy the project's internal logical-consistency check. \\
10 & Penalty or red card & Tournament-wide compound empirical rate & An independence-product of separate penalty and red-card base rates was rejected, since compound "or" questions have a documented history of overstating extremity under an independence assumption. \\
11 & Substitute scores or assists & Base-rate formula validated in the most recent prior match & A per-substitute player-level model was rejected for lack of a stable individual sample for either bench. \\
12 & France 5+ corners & Team-level empirical trend (Poisson at $\hat\lambda \approx 6.3$) & Extrapolating the corner trend's most recent peak (9, against the toughest opponent yet) as if it were the expected value was rejected; the estimate is anchored to the four-game mean, not the peak. \\
13 & Goal after second hydration break & Flat late-window shading, not a fitted lambda-scaling model & A time-fraction-scaled Poisson projection was available and was deliberately not used at face value; this exact question category has the worst documented track record in the campaign. \\
14 & Card in stoppage time & Flat late-window shading & Same rejection as Q13, for a narrower sub-window. \\
15 & France win in regulation & Ordered logit + qualitative Elo cross-check & The Poisson/Dixon--Coles goal model's implied win probability was rejected as the primary source because of its documented systematic compression bias against the stronger side. \\
\bottomrule
\end{longtable}

\section{Results}
\label{sec:results}

Table \ref{tab:finalprices} reports the fifteen final submitted probabilities. No question
was priced outside the range $[0.25, 0.74]$, consistent with the campaign's standing
practice of never submitting at the extreme bounds of the probability simplex.

\begin{table}[H]
\centering
\caption{Final submitted probabilities, France vs. Morocco, 2026-07-09.}
\label{tab:finalprices}
\begin{tabular}{p{0.6cm}p{9.6cm}r}
\toprule
Q & Question & $\hat p$ \\
\midrule
1 & Mbapp\'e (FRA) to score & 0.54 \\
2 & Hakimi (MAR) 1+ shots on target & 0.30 \\
3 & Brahim D\'iaz (MAR) 1+ shots on target & 0.25 \\
4 & Dembel\'e (FRA) 2+ shots on target & 0.30 \\
5 & Morocco 3+ shots on target & 0.60 \\
6 & Total goals $\le 2$ & 0.60 \\
7 & 5+ total cards & 0.32 \\
8 & Tied at halftime & 0.45 \\
9 & Both halves same number of goals & 0.32 \\
10 & Penalty kick awarded OR red card shown & 0.38 \\
11 & Substitute scores or assists & 0.34 \\
12 & France 5+ corner kicks & 0.74 \\
13 & Goal scored after second hydration break & 0.50 \\
14 & Card shown during 1st- or 2nd-half stoppage time & 0.28 \\
15 & France win in regulation & 0.52 \\
\bottomrule
\end{tabular}
\end{table}

Because no crowd or market data exists for this match at the time of writing, no
Relative Brier Points figure can yet be computed for these fifteen submissions; the
settlement record will be appended to this document once the match concludes and the
platform publishes outcomes and crowd consensus values, following the same
per-question audit format used in the campaign's other match-level case studies.

\section{Incorporating Live Market Data}
\label{sec:marketrevision}

The estimates in Table \ref{tab:finalprices} were built with no external market
whatsoever. This section documents a subsequent revision after two live external data
sources were queried: the JTC platform's own web API, to check whether a crowd consensus
had begun to form since the initial pricing pass, and the public Smarkets exchange API,
to obtain a sharp-market anchor of the kind used routinely elsewhere in this campaign.

\subsection{Checking the JTC Crowd Consensus Directly, and a Correction}

Every question record returned by the JTC web API's market endpoint
(\texttt{GET /api/markets}, queried with the event and lobby identifiers) carries a
\texttt{current\_value} field alongside an \texttt{initial\_value} field, both on a
0--100 scale. A direct query of this endpoint, filtered to the fifteen market
identifiers for this fixture, returned \texttt{current\_value} equal to
\texttt{initial\_value} (50) for all fifteen questions. This was initially reported as
confirmation that no net crowd trading activity had occurred on any of the fifteen
questions---that is, as a live crowd-consensus reading of exactly 50\% on every
question.

That reading is incorrect, and the error is worth documenting rather than silently
fixing, since it illustrates a specific and avoidable failure: a newly observed API
field was assumed to mean what its name suggested without being checked against a case
whose answer was already known. The correction was made by pulling every trade recorded
against this event (\texttt{GET /api/trades}) and matching each trade's \texttt{market\_id}
field against the \emph{record} \texttt{id} field of the corresponding market---not, as
first attempted, against that market's own \texttt{market\_id} field, a distinct
identifier that happens to co-exist on the same record and produces zero matches if used
by mistake. This match yielded 209 markets elsewhere in the same event that carry both a
nonzero trade count and a settled status, meaning their true outcomes are independently
known. Across all 209, \texttt{current\_value} takes exactly one of three values, 0, 50,
or 100, with no intermediate readings whatsoever: 50 while a market remains open,
regardless of how many trades have been recorded against it, and 0 or 100 immediately
upon settlement, matching the realized outcome. The field is a binary
open/resolved-no/resolved-yes flag, not a probability, and it is architecturally
incapable of reporting a crowd consensus at any point before a market closes.

This finding does not overturn Section \ref{sec:questionmapping}'s conclusion that no
crowd anchor is available for this match; it removes the basis for having claimed to
verify that conclusion by this particular means. The project's own prior research
(catalogued outside this paper, in a dedicated review of crowd-consensus prediction
literature) had already established that no API access to the platform's pre-close crowd
consensus exists. That finding was correct, is not contradicted by anything reported
here, and should have been the first thing checked before this section's initial claim
was written. \texttt{RULE1}'s crowd-plus-market average remains inapplicable for this
match, for the reason the project's own prior research already gave, not for the reason
first reported.

\subsection{Smarkets Sharp-Market Data}

A query of the public Smarkets event list for World Cup 2026 football fixtures located
event 45178292 (``France vs Morocco'', 2026-07-09), carrying 173 distinct betting
markets. Twelve of these correspond directly or near-directly to questions in this
match's fifteen-question set; contracts and live bid/offer quotes were retrieved for each
and converted to mid prices. Table \ref{tab:marketquotes} reports the key quotes used.

\begin{table}[H]
\centering
\caption{Smarkets mid prices used to revise the fifteen estimates.}
\label{tab:marketquotes}
\begin{tabular}{p{6.4cm}rl}
\toprule
Market & Mid price & Liquidity \\
\midrule
Full-time result: France & 0.608 & Liquid \\
Full-time result: Draw & 0.252 & Liquid \\
Full-time result: Morocco & 0.141 & Liquid \\
Total goals under 2.5 & 0.536 & Liquid \\
Total cards over 4.5 & 0.227 & Liquid (wide spread) \\
Half-time result: Draw & 0.439 & Liquid \\
Penalty to be awarded: No & 0.722 & Liquid (Yes side unquoted; inferred by complementarity) \\
Any player sent off: Yes & 0.147 & Liquid \\
Mbapp\'e anytime goalscorer & 0.503 & Liquid \\
Hakimi 0.5+ shots on target & 0.363 & Liquid \\
Brahim D\'iaz 0.5+ shots on target & 0.435 & Illiquid (offer only) \\
Dembel\'e 1.5+ shots on target & 0.244 & Liquid \\
Morocco 3.5+ shots on target & 0.437 & Liquid \\
France 6.5+ corners & 0.428 & Liquid \\
\bottomrule
\end{tabular}
\end{table}

Two of these quotes require conversion before they can be compared to the fifteen
questions as worded, since Smarkets does not offer a line at exactly the threshold asked.
For Morocco's shots-on-target total, the market's own line sits at 3.5, giving
$P(\text{SOT} \ge 4) = 0.437$; solving for the Poisson rate implied by this figure gives
$\hat\lambda = 3.377$, from which $P(\text{SOT} \ge 3) = 0.656$---a figure that agrees
closely with the pre-market empirical estimate built from Morocco's own competitive-game
average ($\hat\lambda \approx 3.33$), a convergence between an independent market and this
paper's own empirical construction that substantially de-risks Question 5 relative to the
analogous team-count question in a prior match that produced this campaign's worst single
result. The same conversion applied to France's corner line (market at 6.5, giving
$\hat\lambda = 6.212$ and $P(\text{corners} \ge 5) = 0.742$) shows an equally close
agreement with the pre-market empirical estimate of 0.74. For the compound
penalty-or-red-card question, the two Smarkets markets are combined under an
independence assumption, $P(\text{pen or red}) = 1 - (1 - 0.278)(1 - 0.147) = 0.384$,
again closely matching the pre-market tournament-wide base rate of 0.38.

No Smarkets market was found, after checking the full 173-market list, for four of the
fifteen questions: both halves recording the same number of goals, a substitute scoring
or assisting, a goal scored after the second hydration break, and a card shown
specifically during stoppage time. These four retain their pre-market, model- or
base-rate-derived estimates unchanged.

\subsection{Revised Estimates}

Table \ref{tab:revisedprices} reports the fifteen estimates after this revision. The
single largest change is to Question 15: the raw ordered-logit win probability (0.549)
had already been shaded downward to 0.52 for a qualitative adjustment reflecting
Morocco's tournament-long resilience, but the live market places France's win probability
at 0.608, seven percentage points above even the model's unshaded output. Because a
liquid, DIRECT-tier market is this campaign's single most reliable pricing category, and
because the market's number already reflects real money's own assessment of the same
qualitative context that motivated the manual shade, the qualitative adjustment is
retired in favour of the market figure, rounded to 0.60. Player-prop questions with
liquid two-sided Smarkets quotes (Questions 1, 2, 4) were blended toward, rather than
fully replaced by, the market price, consistent with this campaign's standing practice of
retaining a documented personal-history signal alongside a market anchor rather than
discarding it outright. The one illiquid, offer-only quote in the set (Question 3,
Brahim D\'iaz) was deliberately not adopted at face value: the pre-market personal-zero
production evidence for this exact player has already produced the single largest win of
a prior match (Canada versus Morocco, $+44.06$ relative Brier points) against an equally
generous illiquid quote, so the estimate was moved only partway toward the market,
from 0.25 to 0.32, rather than converging on it.

\begin{table}[H]
\centering
\caption{Revised final probabilities after incorporating live Smarkets data.}
\label{tab:revisedprices}
\begin{tabular}{p{0.6cm}p{8.4cm}rrr}
\toprule
Q & Question & Pre-market $\hat p$ & Revised $\hat p$ & Source \\
\midrule
1 & Mbapp\'e (FRA) to score & 0.54 & 0.52 & Blended, liquid \\
2 & Hakimi (MAR) 1+ shots on target & 0.30 & 0.34 & Blended, liquid \\
3 & Brahim D\'iaz (MAR) 1+ shots on target & 0.25 & 0.32 & Partial override, illiquid \\
4 & Dembel\'e (FRA) 2+ shots on target & 0.30 & 0.27 & Blended, liquid \\
5 & Morocco 3+ shots on target & 0.60 & 0.63 & Market-fit lambda \\
6 & Total goals $\le 2$ & 0.60 & 0.57 & Blended, liquid \\
7 & 5+ total cards & 0.32 & 0.26 & Blended, liquid \\
8 & Tied at halftime & 0.45 & 0.44 & Direct, liquid \\
9 & Both halves same number of goals & 0.32 & 0.32 & Unchanged, no market \\
10 & Penalty kick awarded OR red card shown & 0.38 & 0.38 & Combined, corroborated \\
11 & Substitute scores or assists & 0.34 & 0.34 & Unchanged, no market \\
12 & France 5+ corner kicks & 0.74 & 0.74 & Market-fit lambda, corroborated \\
13 & Goal scored after second hydration break & 0.50 & 0.50 & Unchanged, no market \\
14 & Card shown during 1st/2nd-half stoppage time & 0.28 & 0.28 & Unchanged, no market \\
15 & France win in regulation & 0.52 & 0.60 & Direct, liquid \\
\bottomrule
\end{tabular}
\end{table}

\section{Discussion}
\label{sec:discussion}

Two features of this pricing exercise depart from the campaign's typical practice and
merit direct discussion.

First, at the time of the initial pricing pass (Sections \ref{sec:questionmapping}--
\ref{sec:results}), the absence of any market anchor removed the single most reliable
pricing tool in the entire project. Across the four-match knockout-stage audit that
preceded this match, the DIRECT tier---submitting at or near a real market price---beat
the crowd on 73.5\% of thirty-four questions for a net gain of 142.39 relative Brier
points, the strongest and most consistent category the campaign has identified. None of
the fifteen questions in this match could draw on that tier at that stage. The two
team-count questions that substitute for it in the absence of a market, Morocco's
shots-on-target threshold (Question 5) and France's corner threshold (Question 12), are
exactly the question type responsible for the campaign's single worst individual result
to date: a $-80.83$ relative-Brier-point loss on a nearly identical Canada
shots-on-target question, caused by projecting a team's own short-run empirical ceiling
upward without having verified that the team's sample covered a plausible adverse game
state. Both questions were initially priced at or below the empirical floor implied by
each team's performance in its more competitive recent fixtures, rather than blended
upward toward any more optimistic reading of the data, precisely to avoid repeating that
failure mode in a setting with no market price available to soften the consequences of a
miss. Section \ref{sec:marketrevision} subsequently found that a live Smarkets query
independently implies almost the same rate for both statistics (Morocco SOT $\hat\lambda
= 3.377$ against an empirical $\hat\lambda \approx 3.33$; France corners $\hat\lambda =
6.212$ against an empirical estimate of 6.0--6.3), which is the strongest evidence
available that this match's particular application of the empirical-floor approach was
not a repeat of the earlier failure, since a genuinely thin or mis-specified empirical
estimate would have had no reason to agree this closely with an independent market.

Second, the richer StatsBomb data available for this match changes what a plausible but
wrong estimate looks like. In every previous match this campaign priced, a data-poor
question (an illiquid player prop, an obscure derived timing statistic) was recognizable
as fragile because the underlying sample was visibly thin. With ten-to-fourteen-match
StatsBomb histories now available for three of the four named players, an estimate built
purely from that history acquires an unearned appearance of statistical solidity. The
project's own diagnostic test of this exact risk---a live comparison of a pure
StatsBomb-history model against the campaign's actual tournament-aware submission on a
different match---found that the historical model loses in aggregate (93.01 versus
120.30 relative-Brier-point equivalent) despite individually beating the tournament-aware
approach on six of fourteen comparable questions. This paper's estimates therefore treat
every StatsBomb-derived rate, including the specific 2022 France--Morocco semifinal, as a
prior to be weighed against and generally subordinated to the two teams' live
2026 form, rather than as a source of authority in its own right. The clearest instance of
this discipline is Question 1: Mbapp\'e's WC2018/2022 scoring rate (0.857 goals per match)
and his current five-game tournament form (a goal in four of five matches) both argue for
an estimate well above 0.7, but the one StatsBomb-covered meeting between these exact two
sides shows him held to zero shots on target, and Morocco's 2026 defensive record (four
goals conceded in five matches, including clean sheets against Scotland and Canada) is
consistent with the same defensive identity that produced that result in 2022. The final
estimate of 0.54 reflects a genuine, opponent-specific regime check rather than either the
career base rate or the current hot streak taken alone.

A further limitation, distinct from data quality, is that this paper's match-winner and
several prop estimates incorporate a small, explicitly qualitative downward adjustment to
France's raw model-implied win probability, on the stated grounds of Morocco's demonstrated
resilience in high-stakes knockout matches this tournament. This adjustment was kept to two
to three percentage points precisely because the campaign's own research on deliberate
probability-shading (\texttt{STRATEGIC\_MARGIN\_PUSH\_RESEARCH.md}) shows that the Relative
Brier Points scoring rule is strictly proper: any deliberate departure from one's honestly
calibrated belief carries a guaranteed expected cost proportional to the square of the
departure, with no rational justification outside of a narrow, discrete leaderboard-cutoff
scenario that does not apply to this campaign's current standing. The adjustment made here
is better understood as a revision to the calibrated belief itself, informed by qualitative
evidence the quantitative model does not directly encode, rather than as an extremizing
strategy, but the paper flags it explicitly so that any future audit of this match's
settled result can attribute error correctly between the model and this adjustment.

\section{Conclusion}

This match required pricing fifteen questions with no crowd or market anchor whatsoever,
a first for this campaign, using a substantially expanded historical data surface whose
main risk was overapplication rather than scarcity. The paper's methodology addressed both
conditions by (i) re-deriving the campaign's core statistical pipeline---Elo, Poisson/
Dixon--Coles/Negative-Binomial goal modeling, and ordered logit outcome modeling---directly
from this project's own historical panel rather than from any live market, and (ii)
explicitly mapping each question to a single primary tool while recording, for every
question, the plausible alternative tool that was considered and rejected, together with
the reason. The resulting fifteen estimates remain within the campaign's standard
$[0.05, 0.95]$ bounds, apply the two highest-confidence validated patterns from the
campaign's accumulated record (personal-zero-production suppression for Brahim D\'iaz and
Hakimi; refusal to project a team shot or corner ceiling upward from a thin sample absent
a market), and explicitly avoid the two lowest-confidence patterns the campaign has
identified (trusting a raw late-window timing model; treating deliberate probability
extremizing as free upside). A subsequent live check (Section \ref{sec:marketrevision})
found that the JTC crowd consensus is not observable pre-close by any means available to
this project, correcting an initial misreading of an unrelated settlement-status field,
and supplied a public Smarkets sharp-market anchor for twelve of the fifteen questions
instead, most importantly revising the match-winner estimate from 0.52 to a
market-anchored 0.60 and independently corroborating, via closely agreeing implied
Poisson rates, the two team-count estimates
that carried this match's largest single-question tail risk in the absence of a market.
Settlement data will be appended to this document once the match concludes.

\bibliographystyle{apalike}
\bibliography{references}

\end{document}
